\section{Zadanie 7 (Aleksander Wieczorek)}
Niech $A_0,A_1,\dots$ będzie nieskończonym ciągiem zdarzeń i niech $P(\Omega-A_n)\le \frac{1}{3^n}$ dla każdego $n$. Uzasadnij, że
$$P(\bigcap_{n=1}^\infty A_n) \ge \frac{1}{2}$$
\subsection*{Rozwiązanie}
Na początku musimy zauważyć, że:
$$\Omega-\bigcap_{n=1}^\infty A_n=\bigcup_{n=1}^\infty (\Omega-A_n)$$
Najprościej jest to narysować i "to widać". Machając rękami możemy powiedzieć, że prawa strona dla każdego i bierze wszystko co nie leży w $A_i$.
Po zsumowaniu okaże się, że dostaniemy wszystko oprócz przecięcia, czyli lewą stronę.\\
Dalsza część jest już dosyć prosta.\\
Korzystając z powyższej obserwacji dostaniemy: 
$$P(\Omega-\bigcap_{n=1}^\infty A_n)=P(\bigcup_{n=1}^\infty (\Omega-A_n))$$
co daje:
$$1-P(\bigcap_{n=1}^\infty A_n)=P(\bigcup_{n=1}^\infty (\Omega-A_n))$$
$$1-P(\bigcup_{n=1}^\infty (\Omega-A_n))=P(\bigcap_{n=1}^\infty A_n)$$
Teraz przekształcamy tezę:
$$P(\bigcap_{n=1}^\infty A_n) \ge \frac{1}{2}$$
$$1-P(\bigcup_{n=1}^\infty (\Omega-A_n))\ge \frac{1}{2}$$
$$P(\bigcup_{n=1}^\infty (\Omega-A_n))\le \frac{1}{2}$$
Teraz wystarczy tylko pokazać, że $P(\bigcup_{n=1}^\infty (\Omega-A_n))\le \frac{1}{2}$. Z zadania 3. wiemy, że 
$$P(\bigcup_{n=0}^\infty A_n) \le \sum_{n=0}^\infty P(A_n)$$
teraz obliczmy naszą sumę korzystając z $P(\Omega-A_n)\le \frac{1}{3^n}$:
$$\sum_{n=0}^\infty P(\Omega - A_n)\le \frac{1}{3}+\frac{1}{9}+\dots = \frac{\frac{1}{3}}{1-\frac{1}{3}}=\frac{1}{2}$$
Co daje:
$$P(\bigcup_{n=0}^\infty (\Omega-A_n)) \le \sum_{n=0}^\infty P(A_n)\le \frac{1}{2}$$